\documentclass[12pt, letterpaper]{article}
\usepackage[utf8]{inputenc}
\usepackage[T1]{fontenc}
\usepackage[spanish]{babel}
\usepackage{amsmath}
\usepackage[margin=2.5cm]{geometry}

\begin{document}

{\noindent\large\bfseries Emisión de paquetes de $n$-fonones en una molécula excitónica mediante procesos de Stokes}\\[2pt]
{\noindent Jhon Sebastian García Barrera}

\bigskip

Este trabajo estudia la generación de paquetes de $n$ fonones en una molécula excitónica formada
por dos puntos cuánticos semiconductores idénticos acoplados coherentemente mediante interacción de
Förster de constante $J$ a un modo fonónico confinado de una nanocavidad acústica de THz. El
sistema extiende el esquema propuesto por Bin et al.\ (PRL 124, 053601, 2020) para un único punto
cuántico al caso de dos emisores colectivos, determinando cómo la hibridación excitónica modifica
las resonancias de Stokes, las oscilaciones super-Rabi y la pureza estadística de los paquetes
emitidos, así como las ventajas que ofrece frente al emisor único en cuanto a los requisitos
experimentales sobre la calidad de la cavidad.

La diagonalización del Hamiltoniano de Förster en el subespacio de una excitación electrónica
produce dos modos colectivos: el estado brillante $|\Psi_+\rangle$, que acopla al láser con momento
dipolar amplificado por $\sqrt{2}$ respecto al de un emisor individual, y el estado oscuro
$|\Psi_-\rangle$, desacoplado de la dinámica de emisión en el límite de emisores idénticos, que
adquiere población únicamente a través del desfase puro. La dinámica disipativa se modela con la
ecuación maestra de Lindblad (GKLS), incorporando el decaimiento fonónico $\kappa$, el decaimiento
excitónico $\gamma$ y el desfase puro $\gamma_\varphi$ en el régimen
$\gamma, \gamma_\varphi \ll \kappa \ll \omega_b$. El acoplamiento de Förster introduce dos efectos
estructurales: un corrimiento rígido $-J$ en las resonancias de Stokes,
$\Delta_n = -n\omega_b - J$, y un factor superradiante $\sqrt{2}$ en la frecuencia super-Rabi
efectiva, que escala como $\sqrt{2}\,\Omega\,(\lambda/\omega_b)^n/\sqrt{n!}$. Estas expresiones se
derivan analíticamente en tres regímenes de parámetros ---acoplamiento débil, acoplamiento fuerte
electrón-fonón y conducción fuerte--- mediante la partición de Löwdin, con excelente acuerdo frente
a la solución numérica exacta en cada caso.

La solución numérica en QuTiP confirma que las resonancias de $g^{(n)}(0)$ se desplazan en $-J$
respecto al caso de un único emisor, en acuerdo con la predicción analítica, y se mantienen robustas
en un amplio rango de $\lambda$ y $\Omega$. La función $g_2^{(2)}(\tau)$ muestra que la estadística
es sintonizable entre los regímenes \textit{n-phonon gun} (sub-Poissoniano), \textit{n-phonon laser}
(Poissoniano) y estado térmico de paquetes, ajustando la relación $\gamma/\kappa$. Las trayectorias
de Monte Carlo confirman el carácter correlacionado de la cascada: en la resonancia de dos fonones,
los dos \textit{clicks} se producen en una ventana $\Delta\tau \approx 0{,}08\,\kappa^{-1}$, con el
estado oscuro prácticamente desocupado y el sistema retornando al estado base mediante la emisión de
un fotón heraldo detectable en un canal independiente.

Los mapas de pureza $\Pi_n$ para $n = 2, 3, 4$ muestran que para $\lambda/\omega_b \approx 0{,}22$
se obtienen $\Pi_2 \approx 98{,}6\%$ y $\Pi_3 \approx 98{,}9\%$, con $\Pi_3 > 95\%$ sostenido
hasta $Q \sim 90$--$130$, frente a $Q \sim 500$ requeridos por el caso de un único punto cuántico:
una relajación de un factor aproximado de cinco en el requisito sobre la cavidad. Esta ventaja
proviene directamente del factor superradiante, que acelera la cascada de emisión y reduce la
ventana temporal disponible para que los procesos fuera de resonancia contaminen el paquete, haciendo
al sistema inherentemente más robusto frente a pérdidas moderadas. El desfase puro resulta negligible
para valores realistas de $\gamma_\varphi$, aunque el método de Muñoz introduce un error sistemático
creciente que trayectorias de Monte Carlo exactas permitirían superar. Los parámetros empleados son
consistentes con sistemas experimentales de puntos cuánticos en nanocavidades acústicas de THz,
situando las predicciones dentro del alcance de la tecnología actual.

\end{document}